
\chapter{Requirements, Impacts and Constraints}
\label{ch:3}
This chapter states what the system was required to do, what it costs, whom it can affect and under
what constraints it was built. The requirements were fixed before the grid was run and are checked
against measurement in Chapter~\ref{ch:5}; the impacts and constraints are the conditions under which
a detector of this kind may responsibly be deployed.

\section{Final Specifications and Requirements}
\label{sec:req}

\subsection{Functional Requirements}
\label{sec:req.func}
\begin{enumerate}[label=F\arabic*., leftmargin=3.1em]
  \item \textbf{Classify.} Assign every input post one of the corpus's six classes, four of which name
    the group a post attacks and two of which hold abuse that names no group and text that is not
    abusive at all.
  \item \textbf{Score, not only decide.} Expose a probability for each class, so that a deployment can
    choose its own operating point and so that discrimination can be measured without a threshold
    (Section~\ref{sec:4.6}).
  \item \textbf{Run without its teachers.} Perform inference from the student alone. The projections,
    the auxiliary irony head and every teacher are discarded after training
    (Section~\ref{sec:3.7}).
  \item \textbf{Train from unlabelled text.} Accept unlabelled rows alongside the labelled split and
    learn from the teacher's soft labels on them, with the supervised term averaged over the labelled
    rows only so that adding unlabelled text does not dilute supervision (Section~\ref{sec:3.5}).
  \item \textbf{Rebuild from source.} Regenerate every corpus, split, probe and transfer set from
    public sources through the build scripts, and verify a rebuild row by row against a manifest of
    split, label, source and the SHA-1 of the normalised text (Section~\ref{sec:4.2}).
  \item \textbf{Report its own evidence.} Generate every table and interval from the run outputs by
    script rather than by transcription, and score each pre-registered prediction against the
    generated tables (Appendix~\ref{app:predictions}).
\end{enumerate}

\subsection{Non-Functional Requirements}
\label{sec:req.nonfunc}
Table~\ref{tab:requirements} states each non-functional requirement, the target set before the runs,
and what was measured against it.

\begin{table}[htbp]
\centering
\small
\caption[Non-functional requirements and what was measured]{Non-functional requirements, the target
fixed before the grid was run, and the measurement that meets or misses it.}
\label{tab:requirements}
\begin{adjustbox}{max width=\textwidth}
\begin{tabular}{p{2.6cm}p{4.3cm}p{4.9cm}p{1.5cm}}
\toprule
Requirement & Target & Measured & Met \\
\midrule
Latency & Single-post inference in a few milliseconds on ordinary hardware & BERT-mini 3.47 ms at batch
1 and 0.37 ms at batch 32; the BiLSTM 2.57 ms and 0.18 ms (Section~\ref{sec:5.11}) & yes \\
Memory & A deployable model well under the 1,279 MB of the largest teacher & 43 MB for BERT-mini in
fp32 and 33 MB after INT8 quantisation; 255 MB and 132 MB for DistilBERT & yes \\
Compute at inference & The recipe must not add to inference cost & The transfer set is used once, in
training; the deployed student is unchanged (Section~\ref{sec:5.11}) & yes \\
Accuracy & A compact student within reach of the largest teacher & DistilBERT 0.8975 against
BERT-large's 0.8973 macro-F1, at a fifth of the parameters (Section~\ref{sec:5.2}) & yes \\
Resolution & Every claimed difference larger than the test set can resolve & A paired interval on
4,326 posts is about 0.014 wide, so 0.007 is the floor; every in-sample effect falls below it and is
reported as undetectable (Section~\ref{sec:4.7}) & partly \\
Reproducibility & Fixed seeds, one environment, generated tables & Three seeds for every main
configuration, 179 runs from one environment, all tables generated
(Section~\ref{sec:4.8}) & yes \\
Reading of implication & A measurable improvement in telling implied abuse from harmless sarcasm &
No intervention in the grid raises the threshold-free AUC above its in-sample band of 0.756 to 0.777
(Section~\ref{sec:5.10}) & no \\
\bottomrule
\end{tabular}
\end{adjustbox}
\end{table}

The last row is the requirement the work does not meet, and Section~\ref{sec:5.10} reports it as such
rather than as an open question.

\section{Societal Impact}
\label{sec:impact.social}
Automated moderation acts on people, and it fails unevenly. The literature documents systems that
misclassify dialect, reclaimed slurs and in-group speech as abuse
\cite{davidson2017automated, borkan2019nuanced, vidgen2020directions}, and the benign-sarcasm probe
used here measures one such failure directly: a compact student trained on this corpus calls between
14 and 59 per cent of author-labelled sarcastic tweets abusive, depending on the model, while those
tweets attack nobody (Section~\ref{sec:5.10}). A detector deployed at platform scale turns that rate
into a volume of wrongly flagged posts.

The recipe this thesis recommends changes that rate, and not only for the better. As the transfer set
grows, the student fires less readily on sarcasm of both kinds: the false-positive rate on harmless
sarcasm falls from 0.35 to 0.25 on BERT-mini and to 0.18 on the BiLSTM, and recall on ironic abuse
falls with it, from 0.74 to 0.57 and to 0.50. Fewer people are wrongly flagged and more implied abuse
is missed. That is a trade between two harms rather than an improvement, the discrimination behind it
does not change, and Section~\ref{sec:5.10} reports both sides of it so a deployment can choose its
operating point knowingly rather than by accident.

\section{Environmental Impact}
\label{sec:impact.env}
The energy cost of a detector is dominated by inference rather than training, because a moderation
pipeline runs the model on every post while training happens once
\cite{schwartz2020green, patterson2021carbon}. That is the case for compact students, and it is the
case this thesis's measurements support: BERT-mini uses 0.87 GFLOPs per sequence against BERT-large's
78.92, a factor of 91, at 0.0018 of the teacher's macro-F1 deficit once the transfer set is used.

Training for the whole thesis took 18.4 GPU-hours on Kaggle T4 accelerators across five sessions that
ran to completion (Section~\ref{sec:4.8}). The teachers, the implicit specialist and the pre-trained
student were each trained once and resumed by every later session rather than retrained. The
out-of-sample recipe adds one forward pass of one teacher over the unlabelled text and a longer
student run; it adds nothing at inference, which is where the energy is spent. The two failed sessions
are counted in the record even though they produced no result, and the driver now kills any command
that writes nothing for two hours rather than burning a session on a hung process.

\section{Ethical Issues}
\label{sec:ethics}
The corpora contain offensive, hateful and harassing text. They are existing public research datasets,
used under their licences and as text only; no new user data was collected, and no attempt was made to
identify the authors of any post. Where a corpus's labels were not needed, as in the transfer set,
they were discarded before anything read them (Section~\ref{sec:4.4}).

Three ethical constraints bear on what the results can claim. The label for implication is contested:
the two expert annotation efforts that mark it agree on only 48.2 per cent of the hateful texts they
share, so every result on that class is bounded by a disagreement between annotators rather than by
the model (Section~\ref{sec:4.2}). The annotation literature argues that such disagreement is signal
and that a single majority label discards it
\cite{uma2021learning, davani2022dealing, plank2022problem, peterson2019human}; the oracle analysis of
Section~\ref{sec:5.4} finds the committee's unreachable headroom to be made of exactly that
population. And the probes are built by rule from published corpora rather than validated by
annotators, which Section~\ref{sec:7.1} records as a limitation; a 300-item human annotation is in
progress. The models are released for research. Deployment would require human review, calibration to
the platform, and monitoring for disparate impact, none of which this thesis evaluates.

\section{Legal Impact}
\label{sec:legal}
Every corpus is used under its own licence: the cyberbullying tweet corpus as distributed under CC0;
OLID, HatEval, the Davidson et al. corpus and TweetEval's configurations under their research licences;
ISHate under BSL-1.0; the Implicit Hate Corpus and iSarcasmEval for research use. No corpus is
redistributed. The implicit benchmark and the transfer set are released as build scripts together with
a per-row manifest carrying the split, the label, the source and the SHA-1 of the normalised text, so
that a reader can verify a rebuild row by row without either party redistributing text
(Sections~\ref{sec:4.2} and~\ref{sec:4.4}).

The texts are public posts and no author identifiers are retained, which is what keeps the work inside
the terms the corpora were released under. A deployment would sit under different obligations from a
research artefact: platform moderation is increasingly subject to duties of transparency, appeal and
record-keeping, and the threshold-free reporting of Section~\ref{sec:4.6} is the form of evidence
those duties require, because a single recall figure at an unstated threshold cannot support them.

\section{Project Management Plan}
\label{sec:pm}
The work was organised around sessions of bounded length and decisions recorded before the run that
tested them. Compute came in Kaggle sessions of at most twelve hours, so one driver script runs the
whole grid and resumes finished work from the previous session's output; Table~\ref{tab:sessions}
lists the seven sessions, what each added, and the cumulative number of student runs. A CPU smoke test
runs every stage on tiny models before any GPU session begins.

Design decisions are kept in the repository's decision log rather than in the thesis, and the four
that changed the study are recorded with the evidence that forced each one (Section~\ref{sec:6.3}).
The out-of-sample runs were pre-registered: nothing about them was decided after the numbers were
seen, and Section~\ref{sec:5.stats} says what that buys and what it cost when a prediction failed.

\section{Risk Management}
\label{sec:risk}
Table~\ref{tab:risks} lists the risks that materialised and what was done about each. They are not
hypothetical: every row happened.

\begin{table}[htbp]
\centering
\small
\caption[Risks that materialised, and the response to each]{The risks that materialised during the
work, what each would have cost, and the response.}
\label{tab:risks}
\begin{adjustbox}{max width=\textwidth}
\begin{tabular}{p{3.1cm}p{4.6cm}p{5.6cm}}
\toprule
Risk & What it would have cost & Response \\
\midrule
Evaluation contamination & A transfer set that overlaps an evaluation set turns a leak into an
apparent gain & Exact-match screens on normalised text against every split, probe and benchmark file;
one screen alone removed 9,988 rows that occur in the held-out ISHate set
(Section~\ref{sec:4.4}) \\
A benchmark solvable by artefact & Pooling the two implicit-hate corpora gives a benchmark a
classifier can score by recognising the source at 0.91 macro-F1 & Splits built from one corpus, the
other held out whole as an out-of-domain test set (Section~\ref{sec:4.2}) \\
Environment divergence & Identical code, data and seeds scored four points apart on a laptop and on
Kaggle & Every table comes from one environment; comparisons are made only within it and numbers from
elsewhere are excluded rather than reconciled (Section~\ref{sec:7.1}) \\
Lost sessions & Two of seven Kaggle sessions failed, one at teacher caching and one frozen in a
diagnostic cell & Resume from the previous session's output, and a watchdog that kills any command
writing nothing for two hours; both sessions are reported rather than hidden
(Section~\ref{sec:4.8}) \\
Confirmation bias & A grid this large offers many arms to report selectively & Predictions and
decision rules written before each out-of-sample run and scored by script; three of eleven failed and
are reported as failures (Appendix~\ref{app:predictions}) \\
Underpowered claims & Differences smaller than the test set can resolve read as results & The minimum
detectable difference, about 0.007 macro-F1, computed before the results and stated wherever a
difference falls below it (Section~\ref{sec:4.7}) \\
\bottomrule
\end{tabular}
\end{adjustbox}
\end{table}

\section{Economic Analysis}
\label{sec:econ}
The economics of this task are a question about inference, not training. Training the entire grid,
179 student runs and every teacher, took 18.4 GPU-hours; at the rate a commercial provider charges for
a T4-class accelerator that is roughly the cost of a single working day of one engineer, and it was
paid here in free Kaggle quota. The marginal cost of the recipe this thesis recommends over ordinary
fine-tuning is one forward pass of one teacher over 168,000 unlabelled tweets, which is cached once
and reused by every student, plus a student run over a training set about six times larger.

Against that stands what the compact student saves at inference, and it recurs on every post. BERT-mini
answers a single post in 3.47 ms and a batch of 32 in 0.37 ms per post, against 26.56 ms and 21.65 ms
for BERT-large: a factor of 59 at batch 32, on the same machine. DistilBERT matches BERT-large's
macro-F1 at a fifth of the parameters and 5.8 times lower batch-32 latency, and INT8 quantisation
halves its stored size for 0.003 macro-F1. A platform running a detector on every post is therefore
choosing between one large model and roughly sixty compact ones for the same serving budget, and the
measurements here say the compact one need not be materially worse. The transfer set is what buys
that, it is unlabelled, and unlabelled text from the platform is the cheapest input in the pipeline:
no annotation is required, and Section~\ref{sec:5.8} shows that ordinary tweets work nearly as well as
abuse-related ones, so it need not even be curated.
